\question[2]{بخش‌پذیری و هم‌نهشتی}

\subsection{رابطه‌های عاد کردن}
می‌گوییم $a$، $b$ را می‌شمرد اگر $a \divs b$ و در غیر این‌صورت $a \ndivs b$.
برای نمونه $3 \divs 12$ ولی $5 \ndivs 12$. نام این نمادها \lr{\texttt{\textbackslash divs}} و \lr{\texttt{\textbackslash ndivs}} است.

برای اعداد طبیعی $m$ و $n$، کوچک‌ترین مضرب مشترک آن‌ها $\lcm(m,n)$ است.
همچنین توابع رشد مانند چندجمله‌ای $\poly(n)$ و لگاریتم چندجمله‌ای $\polylog(n)$ را داریم.

\subsection{هم‌نهشتی به پیمانه}
وقتی دو عدد تفاضلشان بر $m$ بخش‌پذیر باشد، می‌نویسیم
\[
	a \iequiv{m} b
\]
و می‌خوانیم: $a$ به پیمانه‌ی $m$ با $b$ هم‌نهشت است. این رابطه را به‌صورت $a \imod{m}$ نیز نشان می‌دهیم (با پرانتز و توضیح فارسی).

مثال: $17 \iequiv{5} 2$ زیرا $17-2 = 15$ و $5 \divs 15$. همچنین $17 \imod{5}$.

\subsection{خواص و نامساوی‌ها}
با تعریف جدید نابرابری‌ها، $\leq$ واقعاً $\leqslant$ و $\geq$ واقعاً $\geqslant$ شده‌است.
برای مثال داریم $\floor{x} \leq x \leq \ceil{x}$.

\subsection{یک جدول نمونه}\footnote{این جدول صرفاً جهت نمایش است.}
\begin{center}
	\begin{tabular}{|c|c|c|}
		\hline
		\lr{Symbol} & نام       & کاربرد         \\
		\hline
		$\divs$     & عاد کردن  & بخش‌پذیری       \\
		$\lcm$      & ک.م.م     & مضرب مشترک     \\
		$\poly(n)$  & چندجمله‌ای & تحلیل الگوریتم \\
		\hline
	\end{tabular}
\end{center}